\section{Sonuc}

Bu projede, EGO-benzeri bir otobus takip sisteminin sentetik veri ile modellenmesi ve bu verinin AWS tabanli gercek zamanli bir akis mimarisi uzerinde islenmesi basarili bicimde gerceklestirilmistir. Sistem; veri uretimi, cihaz-bulut haberlesmesi, akisa dayali isleme, depolama ve gorsellestirme katmanlarini butunlesik bir sekilde bir araya getirmektedir. Ham veri ile turetilmis veri arasindaki ayrimin korunmasi ve karar ureten alanlarin bulutta hesaplanmasi, calismanin teknik anlamini guclendirmistir.

Sonuc olarak, bu calisma bulut bilisim dersinin temel hedefleriyle uyumlu bir sekilde, MQTT tabanli gercek zamanli veri akisi ile AWS servis entegrasyonunu pratikte gosteren bir referans uygulama ortaya koymustur. Proje ayni zamanda, sentetik veri kullanilsa dahi dogru mimari kararlar, uygun servis secimi ve sistematik dogrulama adimlari ile akademik olarak savunulabilir, tekrar edilebilir ve gosterilebilir bir sistem kurulabilecegini gostermistir.
